\section{Methodology}
\label{sec:methods}

\subsection{Governing physical dynamics}
The ambient concentration field $C_i(\mathbf{x},t)$ of chemical species $i$ evolves
according to the atmospheric advection--diffusion--reaction equation,
\begin{equation}
  \frac{\partial C_i}{\partial t}
  = \underbrace{-\,\nabla \cdot (\mathbf{u}\, C_i)}_{\text{advection}}
  + \underbrace{\nabla \cdot (\mathbf{K}\, \nabla C_i)}_{\text{turbulent diffusion}}
  + \underbrace{R_i(C_1,\dots,C_m,\,T,\,h\nu)}_{\text{chemistry}}
  + \underbrace{E_i(\mathbf{x},t)}_{\text{emissions}},
  \label{eq:adr}
\end{equation}
where $\mathbf{u}(\mathbf{x},t)$ is the wind velocity field, $\mathbf{K}$ the
(second-order) eddy-diffusivity tensor, $R_i$ the net chemical production/loss rate for
species $i$ as a function of all $m$ reacting species, temperature $T$ and the actinic
photon flux $h\nu$, and $E_i$ the surface emission rate. Equation~\eqref{eq:adr} makes
the epistemic gap explicit: a single-site, contemporaneous regression of $C_i$ on the
other $C_j$ and a few meteorological scalars discards the spatial transport operators
entirely, linearizes the strongly non-linear reaction term $R_i$, and---critically---
never observes the emission field $E_i$ on which policy acts. It can estimate a local,
time-averaged \emph{response surface}, not the operator in \eqref{eq:adr}.

\subsection{Predictive formulation and leakage controls}
Let $y_{s,t}$ denote the standardized ozone concentration at station $s$ and hour $t$,
and $\mathbf{x}_{s,t}\in\RR^{d}$ the vector of contemporaneous covariates (the other
six pollutants and four meteorological variables). We estimate $\hat y = f(\mathbf{x})$
with gradient-boosted regression trees \citep{chen2016xgboost}, chosen for their strong
tabular performance under non-linear interactions and implemented within a distributed
Apache Spark pipeline \citep{zaharia2016spark}; hyper-parameters are fixed a priori and
never tuned on any evaluation fold.

Because $y_{s,t}$ is strongly autocorrelated in both time and space, the choice of
evaluation protocol is not a detail but the central inferential control
\citep{kaufman2012leakage,roberts2017cv,varma2006cv}. We contrast three protocols:
\begin{enumerate}
  \item \emph{i.i.d.\ random split} (70/30 over rows): the standard but
  \emph{leaky} baseline, in which near-duplicate neighbouring hours of the same station
  appear in both train and test.
  \item \emph{Blocked temporal holdout}: train on quarters Q1--Q3 (months 1--9),
  evaluate on Q4 (months 10--12). Temporal blocks are never split.
  \item \emph{Spatial leave-stations-out}: $5$-fold \texttt{GroupKFold} with the station
  identifier as the grouping key, so that no station appears in both train and test;
  this probes transfer to \emph{unmonitored} locations.
\end{enumerate}
All preprocessing (standardization, imputation) is fit on the training fold only.
We report against a ladder of reference baselines: (i) \emph{lag-1 persistence},
$\hat y_{s,t}=y_{s,t-1}$; (ii) a \emph{station $\times$ hour $\times$ month
diurnal-seasonal mean}, the strongest purely climatological predictor; and (iii) an
$L_2$-regularized \emph{Ridge} regression \citep{hastie2009esl} on the same covariates,
as a linear reference for the boosted model. Uncertainty is quantified by a station-block bootstrap.

\subsection{Covariance clustering versus bilinear receptor modelling}
The pipeline's first stage groups species by $k$-means over their inter-species
Pearson correlation profiles. It is essential to state precisely what this does
\emph{not} do. Quantitative source apportionment is a bilinear receptor model; positive
matrix factorization (PMF) \citep{paatero1994pmf} factorizes a speciated
concentration matrix $X=[x_{ij}]$ (sample $i$, chemical species $j$) as
\begin{equation}
  x_{ij} = \sum_{k=1}^{p} g_{ik} f_{kj} + e_{ij},
  \qquad g_{ik}\ge 0,\; f_{kj}\ge 0,
  \label{eq:pmf}
\end{equation}
where $g_{ik}$ is the contribution of source $k$ to sample $i$, $f_{kj}$ the
chemical profile of source $k$, and $e_{ij}$ a residual weighted by the measurement
uncertainty $\sigma_{ij}$ through the objective
$Q=\sum_{ij}(e_{ij}/\sigma_{ij})^2$. PMF requires (a) chemically speciated composition
(including organic and elemental carbon and key tracers), (b) per-measurement
uncertainties, and (c) enough species to identify $p$ non-negative factors. Our data
provide only seven criteria pollutants, no speciation, no VOCs, and no per-measurement
uncertainties. $k$-means over criteria-gas correlations is therefore an
\emph{exploratory inter-species covariance clustering}, and its clusters are
mathematically distinct from the latent source factors $\{g_{ik},f_{kj}\}$ of
\eqref{eq:pmf}; they admit no mass-balance or source interpretation.

\subsection{Causal discovery via DirectLiNGAM}
For exploratory structure, we model the standardized pollutant vector
$\zvec\in\RR^{m}$ as a linear, non-Gaussian, acyclic structural equation model,
\begin{equation}
  \zvec = \bmat\,\zvec + \evec,
  \label{eq:sem}
\end{equation}
where $\bmat$ is a strictly lower-triangular (under some causal ordering) matrix of
structural coefficients and $\evec$ is a vector of mutually independent, non-Gaussian
exogenous disturbances. Unlike constraint-based
discovery \citep{spirtes2000causation}, DirectLiNGAM
\citep{shimizu2006lingam,shimizu2011directlingam} estimates a causal ordering and the
coefficients $\bmat$ directly, exploiting non-Gaussianity to orient edges; non-linear
generalizations exist via additive-noise models \citep{hoyer2009anm} but inherit the
same confounding sensitivity. Its validity rests on four assumptions:
(A1) \emph{acyclicity} of the generating graph; (A2) \emph{linearity} of the structural
equations; (A3) \emph{non-Gaussian} disturbances with non-zero variance; and
(A4) \emph{causal sufficiency}---no unobserved common causes of two or more observed
variables. We encode domain knowledge as a prior-knowledge constraint matrix
$K\in\{0,1,{-}1\}^{m\times m}$ that forbids edges from secondary pollutants
(e.g.\ O\textsubscript{3}) into primary combustion tracers (e.g.\ NO, CO), consistent
with emission chemistry. Edge reliability is assessed by a station-block bootstrap
(reporting per-edge selection frequency and coefficient sign), and---decisively---by
refitting with the meteorological variables included as exogenous nodes, which tests
assumption (A4) against the dominant confounder.

\subsection{Analytical derivation of the global Pareto frontier}
The optimization stage searches, with the NSGA-II evolutionary algorithm
\citep{deb2002nsga2} as implemented in the Optuna framework \citep{akiba2019optuna},
for policies minimizing two standardized objectives built from the fitted linear
response surface. Empirically the fitted surface is
\begin{align}
  z_{\mathrm{O_3}}   &= \alpha\, z_{\mathrm{NO}}, & \alpha &< 0, \label{eq:resp_o3}\\
  z_{\mathrm{PM_{2.5}}} &= \beta_1 z_{\mathrm{NO}} + \beta_2 z_{\mathrm{NO_2}}
                            + \beta_3 z_{\mathrm{CO}}, & \beta_1,\beta_2,\beta_3 &> 0,
  \label{eq:resp_pm}
\end{align}
with fitted values $\alpha=-0.124$, $\beta_1=0.092$, $\beta_2=0.219$, $\beta_3=0.228$;
SO\textsubscript{2} and VOC enter with zero coefficient (VOC is not a fitted variable),
so the effective decision vector is
$\mathbf{z}=(z_{\mathrm{NO}},z_{\mathrm{NO_2}},z_{\mathrm{CO}})$ over the search
hypercube $z_j\in[z_{\min},z_{\max}]=[-2,0]$ (reductions only, bounded at $-2\sigma$).

\begin{lemma}[Collapse of the linear Pareto front]
\label{lem:pareto}
Consider the bi-objective minimization of
$\big(f_1(\mathbf{z}),f_2(\mathbf{z})\big)=\big(z_{\mathrm{O_3}},z_{\mathrm{PM_{2.5}}}\big)$
defined by \eqref{eq:resp_o3}--\eqref{eq:resp_pm} over
$\mathbf{z}\in[z_{\min},z_{\max}]^3$ with $\alpha<0$ and $\beta_1,\beta_2,\beta_3>0$.
Then every Pareto-optimal point satisfies
$z_{\mathrm{NO_2}}^{\star}=z_{\mathrm{CO}}^{\star}=z_{\min}$, and the Pareto set is the
one-dimensional segment
\[
  \mathcal{P}=\Big\{\big(z_{\mathrm{NO}},z_{\min},z_{\min}\big)\;:\;
  z_{\mathrm{NO}}\in[z_{\min},z_{\max}]\Big\},
\]
whose image is the line segment
$z_{\mathrm{PM_{2.5}}} = \tfrac{\beta_1}{\alpha}\,z_{\mathrm{O_3}}
+(\beta_2+\beta_3)\,z_{\min}$.
\end{lemma}

\begin{proof}
The objective $f_1$ depends only on $z_{\mathrm{NO}}$, while $f_2$ is strictly
increasing in each of $z_{\mathrm{NO_2}}$ and $z_{\mathrm{CO}}$ (as $\beta_2,\beta_3>0$)
and does not appear in $f_1$. Fix any feasible point $\mathbf{z}$ with
$z_{\mathrm{NO_2}}>z_{\min}$. The point $\mathbf{z}'$ obtained by lowering
$z_{\mathrm{NO_2}}$ to $z_{\min}$ (all else equal) is feasible, leaves $f_1$ unchanged,
and strictly decreases $f_2$; hence $\mathbf{z}'$ dominates $\mathbf{z}$, so $\mathbf{z}$
is not Pareto-optimal. The identical argument applies to $z_{\mathrm{CO}}$. Therefore
every Pareto-optimal point has $z_{\mathrm{NO_2}}^{\star}=z_{\mathrm{CO}}^{\star}=z_{\min}$.

Restricted to this face the objectives reduce to
$f_1=\alpha z_{\mathrm{NO}}$ and
$f_2=\beta_1 z_{\mathrm{NO}} + (\beta_2+\beta_3)z_{\min}$. Since $\alpha<0$, $f_1$ is
strictly decreasing in $z_{\mathrm{NO}}$, whereas $f_2$ is strictly increasing in
$z_{\mathrm{NO}}$ (as $\beta_1>0$); the two objectives are therefore in strict conflict
along $z_{\mathrm{NO}}$, so no interior point of the segment dominates another and the
entire segment is Pareto-optimal. Eliminating $z_{\mathrm{NO}}=f_1/\alpha$ gives
$f_2=\tfrac{\beta_1}{\alpha}f_1+(\beta_2+\beta_3)z_{\min}$, a straight line.
\end{proof}

\begin{corollary}
\label{cor:pareto}
The true Pareto frontier of the fitted linear surrogate is a closed-form line segment
computable without any metaheuristic search. Any evolutionary approximation can at best
match it and, lacking the boundary structure of \cref{lem:pareto}, will generically
return interior, dominated points.
\end{corollary}

\paragraph{Benchmark metrics.}
Let $\mathcal{F}$ be a dense discretization of the exact frontier $\mathcal{P}$ and
$A=\{\mathbf{a}_k\}_{k=1}^{n}$ the archive of heuristic solutions in objective space.
The \emph{generational distance} is
\begin{equation}
  \mathrm{GD}(A)=\Big(\frac{1}{n}\sum_{k=1}^{n} d_k^{2}\Big)^{1/2},
  \qquad d_k=\min_{\mathbf{p}\in\mathcal{F}}\lVert \mathbf{a}_k-\mathbf{p}\rVert_2 ,
  \label{eq:gd}
\end{equation}
and, for a shared reference point $\mathbf{r}$ dominated by all points, the
\emph{hypervolume error} is
\begin{equation}
  \mathrm{HVE}=\frac{\mathrm{HV}(\mathcal{F};\mathbf{r})-\mathrm{HV}(A;\mathbf{r})}
  {\mathrm{HV}(\mathcal{F};\mathbf{r})},
  \label{eq:hve}
\end{equation}
where $\mathrm{HV}(\cdot\,;\mathbf{r})$ is the dominated hypervolume \citep{zitzler1999hv}.
$\mathrm{GD}=0$ and $\mathrm{HVE}=0$ iff the archive lies on the true frontier.
